The evaluation protocol is designed to prevent a common overclaim: confusing stronger immediate forgetting with durable forgetting.

\paragraph{Immediate matching.}
For each method under test, candidate baselines are swept and selected so that their $r@0$ is within a tolerance band of the reference method.
The local default tolerance is $0.05$.
Candidates are rejected when retain NLL or refusal rate is outside the allowed band.
This creates a fairer comparison: both methods start with approximately the same number of resurrected forget items.

\paragraph{Conditional resurrection.}
Let
\[
  F_0 = \{i \in F : m_i(\theta_0) \le \tau_i\}
\]
be the set of items forgotten at $t=0$.
Conditional resurrection at horizon $T$ is
\[
  c_T = \frac{1}{|F_0|}\sum_{i \in F_0}\mathbf{1}[m_i(\theta_T) > \tau_i].
\]
This directly measures whether forgotten items stay forgotten.

\paragraph{Conditional survival AUC.}
Conditional survival is $1-c_t$.
We compute running AUC over the fixed horizon grid.
The most informative horizons in the current harness are $t=32$, $t=128$, and $t=512$.

\paragraph{Stress pools.}
Benign relearn pools are useful for smoke tests, but answer-bearing pools are decisive for durability.
The harness includes declarative, paraphrased QA, and mixed answer-bearing pools, plus held-out versions with different wording.
Held-out pools prevent tuning only to one narrow relearning phrasing.

\paragraph{Pass condition.}
A positive HLC claim requires lower conditional resurrection or higher conditional survival AUC than the strongest immediate-matched baseline, while retaining similar utility and low refusal.
If HLC only wins on $r@0$, the result is not a durability win.
